\begin{textsample}{POS Dim 3 – rj – Score 30.00 – rj/rj\_inf\_cam\_exp\_16.txt}  \label{ex:f3_pos_001}
Práticas pedagógicas que auxiliam na garantia dos Direitos de Aprendizagem no Campo de Experiências Espaços, Tempos, Quantidades, Relações e Transformações Nas atividades exploratórias das quais participam interagindo com os \textbf{colegas}, com a \textbf{professora} e com o material disponível, as \textbf{crianças} podem se apropriar de formas produtivas de pensar os mundos da natureza e da sociedade, incluindo os animais, as plantas, os objetos, a tecnologia, o comportamento humano e outros aspectos da cultura, observar características, diferenças, regularidades e irregularidades de fenômenos, e procurar explicar os modos destes se constituírem e se transformarem.
Nesse ponto elas vivenciam de modo integrado experiências em relação ao tempo, ao espaço, às quantidades, relações e transformações.
Nesse contexto, a \textbf{professora} assume o papel de mediador das relações das \textbf{crianças} com os conhecimentos já elaborados acerca da natureza e da sociedade, \textbf{cuidando} para que desfrutem e se surpreendam com as descobertas que fazem, alegrem -se com suas próprias capacidades de conhecer e sintam interesse e paixão por essas atividades.
Para tanto ela deve \textbf{acolher} sentimentos, questões e ideias das \textbf{crianças} e propor -lhes perguntas, questionamentos, dúvidas que as mobilizem a indagar acerca de algum aspecto do mundo na construção de novos conhecimentos.
Ela deixa de ser um informante dos conhecimentos tidos como científicos, um transmissor de conteúdos para que as \textbf{crianças} aprendam, e passa a ser um investigador de como elas pensam, na medida em que interpreta suas hipóteses, considera seus argumentos e analisa suas experiências.
O básico é partir das ideias e representações que as \textbf{crianças} possuem e fazer -lhes perguntas instigantes, oferecer -lhes meios para que busquem mais informações e possam reformular suas ideias iniciais, respeitando os raciocínios, relações, comparações e analogias que elas elaboram e fazer -lhes questionamentos que as leve a pensar sobre eles e considerar a enorme quantidade de ideias diferentes que se têm para explicar coisas e fenômenos.
O importante é que tenham tempo suficiente para explorar e possam viver de forma \textbf{repetida} diferentes experiências e oportunidades.
As experiências das quais participam devem propiciar às \textbf{crianças} a exploração de objetos ( observar a água em forma de gelo, a água na forma líquida e o vapor d’água ), a formulação de perguntas ( como a água evaporou? ), a construção de hipóteses ( é porque está calor? ), o desenvolvimento de generalizações ( o sorvete também derrete porque está calor? ) e a aprendizagem de palavras específicas ( evaporar, derreter, líquido, \textbf{sólido} etc. ).
Isso requer \textbf{ouvir} com atenção as conversas e questionamentos das \textbf{crianças} e problematizar os assuntos por elas trazidos, incentivá -las a comunicar umas às outras suas descobertas ( “ \textbf{Conte} aos \textbf{colegas} o que você descobriu sobre a duração da vida das baleias! ” ), suas previsões ( “ o que vocês acham que aconteceria se todos os relógios do mundo fossem quebrados?
O que vai acontecer se eu jogar esse algodão na água?
E se eu jogar essa bolinha de \textbf{plástico}? ” ), hipóteses ( “ Por que vocês acham que a tampa vai boiar? ” ), a descrever suas experiências ( “ O que aconteceu com o cubo de gelo quando vocês jogaram água nele? ” ), bem como registrar suas ideias, observações e investigações ( “ Vou dar a cada um de vocês uma prancheta e um \textbf{lápis} e vocês vão sair pelo \textbf{parque} observando os bichos que encontrarem e registrando como eles são e o que eles fazem ” ).
Também podem ser criadas oportunidades para as \textbf{crianças} \textbf{demonstrarem} o que aprenderam utilizando diferentes linguagens, seja desenhando, fazendo um gráfico, anotando números, escrevendo palavras ou frases simples, ou apresentando oralmente suas investigações ou conclusões aos \textbf{colegas}.
Novamente, a transdisciplinaridade dos campos de experiências assegura a interação de diferentes linguagens na apropriação do mundo pela \textbf{criança}.
Este processo é gradativo e muito dependente das oportunidades criadas nas unidades de Educação \textbf{Infantil} urbanas e do campo para todas as \textbf{crianças}.
A BNCC ( BRASIL, 2017 ), de modo a orientar os projetos políticos pedagógicos das unidades de Educação \textbf{Infantil}, propôs que neles as \textbf{crianças} tivessem garantidos os seguintes direitos mediadores de significativas aprendizagens que devem se relacionar com os campos de experiências. - CONVIVER com outras \textbf{crianças} e \textbf{adultos}, em \textbf{pequenos} e grandes grupos, utilizando diferentes linguagens, ampliando o conhecimento de si e do outro, o respeito em relação à cultura e às diferenças entre as pessoas. - \textbf{BRINCAR} de diversas formas, em diferentes espaços e tempos, com diferentes \textbf{parceiros} ( \textbf{crianças} e \textbf{adultos} ) de forma a ampliar e diversificar suas possibilidades de acesso a produções culturais.
A participação e as transformações introduzidas pelas \textbf{crianças} nas \textbf{brincadeiras} devem ser valorizadas, tendo em vista o \textbf{estímulo} ao desenvolvimento de seus conhecimentos, sua imaginação, criatividade, experiências emocionais, corporais, \textbf{sensoriais}, expressivas, cognitivas, sociais e \textbf{relacionais}. - PARTICIPAR ativamente, com \textbf{adultos} e outras \textbf{crianças}, tanto do planejamento da gestão da escola e das atividades propostas pelo educador quanto da realização das atividades da vida cotidiana, tais como a escolha das \textbf{brincadeiras}, dos materiais e dos ambientes, desenvolvendo diferentes linguagens e elaborando conhecimentos, decidindo e se posicionando. - EXPLORAR movimentos, \textbf{gestos}, \textbf{sons}, formas, \textbf{texturas}, cores, palavras, emoções, transformações, relacionamentos, histórias, objetos, elementos da natureza, na escola e fora dela, ampliando seus saberes sobre a cultura em suas diversas modalidades : as artes, a escrita, a ciência e a tecnologia. - EXPRESSAR, como sujeito dialógico, criativo e sensível, suas necessidades, emoções, sentimentos, dúvidas, hipóteses, descobertas, opiniões, questionamentos, por meio de diferentes linguagens. - CONHECER -SE e construir sua identidade pessoal, social e cultural, constituindo uma imagem positiva de si e de seus grupos de pertencimento, nas diversas experiências de \textbf{cuidados}, interações, \textbf{brincadeiras} e linguagens vivenciadas na \textbf{instituição} escolar e em seu contexto familiar e comunitário.
À \textbf{criança} devem ser garantidos esses direitos, a fim de que esta alcance sua participação e autonomia não só na escola, mas em toda a sociedade.
Conceder direitos não é uma grande façanha, é apenas o começo de uma revolução que há tempos tornou -se necessária dentro das \textbf{instituições} escolares.
A BNCC reconhece o aumento da complexidade da aprendizagem à medida que as \textbf{crianças} crescem.
Isso mostra a necessidade de estruturação e organização de situações de aprendizagem.
Esses direitos garantem uma concepção de \textbf{criança} como ser observador, questionador, capaz de levantar hipóteses, concluir, julgar e \textbf{assimilar} valores.
Isso contribui para que possa construir seus conhecimentos e apropriar -se deles de forma sistematizada, por meio da ação e nas interações com o mundo físico e social [ e ] não deve resultar no confinamento dessas aprendizagens a um processo de desenvolvimento natural ou espontâneo.
Ao contrário, reitera a importância e necessidade de imprimir intencionalidade educativa às práticas pedagógicas na Educação \textbf{Infantil}, tanto na \textbf{creche} quanto na \textbf{pré-escola} ( BRASIL, 2017, p. 35 ).

% matched lemmas: acolher, adulto, assimilar, brincadeira, brincar, colega, contar, creche, criança, cuidado, cuidar, demonstrar, estímulo, gesto, infantil, instituição, lápis, ouvir, parceiro, parque, pequeno, plástico, professora, pré-escola, relacional, repetir, sensorial, som, sólido, textura
\end{textsample}
